% GENERATED FILE. Edit canonical lessons, then run npm run generate:books.
% canonical-source-hash: fnv1a64:9349969834be0c76
% canonical-lessons: BN-C67-bali, BN-C67-chad, BN-C67-somudro, BN-C67-agun, BN-C67-boroph

\chapter{Sand, Moon, Sea, Fire, Ice}
\label{ch:bn-sand-moon-sea-fire-ice}
\hlchaptermodality{67}

\begin{chapteropening}
\textbf{By the end of this chapter:} \emph{I can say sand, the moon, the sea, fire and ice.}

Complete the last lesson of chapter 67: I can say sand, the moon, the sea, fire and ice.
\end{chapteropening}

\section[\texorpdfstring{\bn{বালি} (bāli)}{বালি (bāli)}]{\bn{বালি} (bāli) — sand}

\label{lesson:BN-C67-bali}



\begin{quote}
Before the new one: say the Bengali for a pond, then the Bengali for a stone.
\end{quote}

\subsection*{You'll want to know: \bn{বালি}}
\textbf{\bn{বালি}} — \emph{bāli} — \textquotedblleft{}sand\textquotedblright{}.

\begin{cousinweb}
\textbf{\bn{বালি}} (\emph{bāli}) is a tadbhava, a word inherited from Sanskrit and reshaped by sound change. It comes from \emph{vālukā}, \textquotedblleft{}sand, gravel\textquotedblright{} (CDIAL 11580). CDIAL notes that its ending \emph{-ukā} was replaced by \emph{-ikā}.
\end{cousinweb}

\begin{grammarlens}[title={What you've built}]
One more word for this chapter.
\end{grammarlens}

\subsection*{Your turn}
\begin{itemize}
\raggedright
  \item \emph{Say it:} \emph{bāli}
  \item \emph{Say it:} \emph{bāli}, once more
  \item \emph{Say it:} say \emph{bāli}
  \item \emph{Recall:} say the Bengali for a pond, then the Bengali for a stone, then say \emph{bāli} again
\end{itemize}

\begin{tcolorbox}[breakable,colback=teal!4,colframe=teal!35!black,title={Before you move on}]
What does \bn{বালি} mean? (\textquotedblleft{}Sand\textquotedblright{}.) Say it once more.
\end{tcolorbox}
\section[\texorpdfstring{\bn{চাঁদ} (chā̃d)}{চাঁদ (chā̃d)}]{\bn{চাঁদ} (chā̃d) — the moon}

\label{lesson:BN-C67-chad}



\begin{quote}
Before the new one: say the Bengali for a stone, then the Bengali for sand.
\end{quote}

\subsection*{You'll want to know: \bn{চাঁদ}}
\textbf{\bn{চাঁদ}} — \emph{chā̃d} — \textquotedblleft{}the moon\textquotedblright{}.

\begin{cousinweb}
\textbf{\bn{চাঁদ}} (\emph{chāṁd}) is a tadbhava, a word inherited from Sanskrit and reshaped by sound change. It comes from \emph{candra}, \textquotedblleft{}shining; moon\textquotedblright{} (CDIAL 4661).
\end{cousinweb}

\begin{grammarlens}[title={What you've built}]
One more word for this chapter.
\end{grammarlens}

\subsection*{Your turn}
\begin{itemize}
\raggedright
  \item \emph{Say it:} \emph{chā̃d}
  \item \emph{Say it:} \emph{chā̃d}, once more
  \item \emph{Say it:} say \emph{chā̃d}
  \item \emph{Recall:} say the Bengali for a stone, then the Bengali for sand, then say \emph{chā̃d} again
\end{itemize}

\begin{tcolorbox}[breakable,colback=teal!4,colframe=teal!35!black,title={Before you move on}]
What does \bn{চাঁদ} mean? (\textquotedblleft{}The moon\textquotedblright{}.) Say it once more.
\end{tcolorbox}
\section[\texorpdfstring{\bn{সমুদ্র} (sômudrô)}{সমুদ্র (sômudrô)}]{\bn{সমুদ্র} (sômudrô) — the sea}

\label{lesson:BN-C67-somudro}



\begin{quote}
Before the new one: say the Bengali for sand, then the Bengali for the moon.
\end{quote}

\subsection*{You'll want to know: \bn{সমুদ্র}}
\textbf{\bn{সমুদ্র}} — \emph{sômudrô} — \textquotedblleft{}the sea\textquotedblright{}.

\begin{cousinweb}
\textbf{\bn{সমুদ্র}} (\emph{sômudrô}) is Sanskrit \emph{samudra}, \textquotedblleft{}sea\textquotedblright{}, taken over whole as a learned loan, a tatsama.
\end{cousinweb}

\begin{grammarlens}[title={What you've built}]
One more word for this chapter.
\end{grammarlens}

\subsection*{Your turn}
\begin{itemize}
\raggedright
  \item \emph{Say it:} \emph{sômudrô}
  \item \emph{Say it:} \emph{sômudrô}, once more
  \item \emph{Say it:} say \emph{sômudrô}
  \item \emph{Recall:} say the Bengali for sand, then the Bengali for the moon, then say \emph{sômudrô} again
\end{itemize}

\begin{tcolorbox}[breakable,colback=teal!4,colframe=teal!35!black,title={Before you move on}]
What does \bn{সমুদ্র} mean? (\textquotedblleft{}The sea\textquotedblright{}.) Say it once more.
\end{tcolorbox}
\section[\texorpdfstring{\bn{আগুন} (āgun)}{আগুন (āgun)}]{\bn{আগুন} (āgun) — fire}

\label{lesson:BN-C67-agun}



\begin{quote}
Before the new one: say the Bengali for the moon, then the Bengali for the sea.
\end{quote}

\subsection*{You'll want to know: \bn{আগুন}}
\textbf{\bn{আগুন}} — \emph{āgun} — \textquotedblleft{}fire\textquotedblright{}.

\begin{cousinweb}
\textbf{\bn{আগুন}} (\emph{āgun}) is a tadbhava, a word inherited from Sanskrit and reshaped by sound change. It comes from \emph{agni}, \textquotedblleft{}fire\textquotedblright{} (CDIAL 14198). CDIAL's addenda list \emph{āgun} and its variants; the main entry gives the shorter \emph{āg} (CDIAL 55).
\end{cousinweb}

\begin{grammarlens}[title={What you've built}]
One more word for this chapter.
\end{grammarlens}

\subsection*{Your turn}
\begin{itemize}
\raggedright
  \item \emph{Say it:} \emph{āgun}
  \item \emph{Say it:} \emph{āgun}, once more
  \item \emph{Say it:} say \emph{āgun}
  \item \emph{Recall:} say the Bengali for the moon, then the Bengali for the sea, then say \emph{āgun} again
\end{itemize}

\begin{tcolorbox}[breakable,colback=teal!4,colframe=teal!35!black,title={Before you move on}]
What does \bn{আগুন} mean? (\textquotedblleft{}Fire\textquotedblright{}.) Say it once more.
\end{tcolorbox}
\section[\texorpdfstring{\bn{বরফ} (bôrôph)}{বরফ (bôrôph)}]{\bn{বরফ} (bôrôph) — ice}

\label{lesson:BN-C67-boroph}



\begin{quote}
Before the new one: say the Bengali for the sea, then the Bengali for fire.
\end{quote}

\subsection*{You'll want to know: \bn{বরফ}}
\textbf{\bn{বরফ}} — \emph{bôrôph} — \textquotedblleft{}ice\textquotedblright{}.

\begin{cousinweb}
\textbf{\bn{বরফ}} (\emph{bôrôph}) is a loan, usually traced to Persian \emph{barf}, \textquotedblleft{}snow, ice\textquotedblright{}.
\end{cousinweb}

\begin{grammarlens}[title={What you've built}]
One more word for this chapter.
\end{grammarlens}

\subsection*{Your turn}
\begin{itemize}
\raggedright
  \item \emph{Say it:} \emph{bôrôph}
  \item \emph{Say it:} \emph{bôrôph}, once more
  \item \emph{Say it:} say \emph{bôrôph}
  \item \emph{Recall:} say the Bengali for the sea, then the Bengali for fire, then say \emph{bôrôph} again
\end{itemize}

\begin{tcolorbox}[breakable,colback=teal!4,colframe=teal!35!black,title={Before you move on}]
What does \bn{বরফ} mean? (\textquotedblleft{}Ice\textquotedblright{}.) Say it once more.
\end{tcolorbox}
